\subsection{共动视界扫描的傅里叶完备性：有限缺陷口径下一维剖面的可逆层析}\label{subsec:cdim-comoving-horizon-scan-fourier-completeness}
设存在有限离临界线零点多重集 \(\{\rho_j\}_{j=1}^{M}\)（按重数计），其中
\(\rho_j=\beta_j+\mathrm{i}\gamma_j\) 且 \(\beta_j<\tfrac12\)。
令 \(\delta_j:=\tfrac12-\beta_j>0\)，并记重数 \(m_j\in\ZZ_{\ge 1}\)。
对任意共动参数 \(x_0\in\RR\)，定义移位图表下的圆盘像
\[
w_{\rho_j,x_0}:=\mathcal{C}\!\bigl((\gamma_j-x_0)+\mathrm{i}\delta_j\bigr)\in\mathbb{D},
\qquad
h_{x_0}(\rho_j):=1-|w_{\rho_j,x_0}|^2,
\]
并定义聚合剖面函数
\[
H(x_0):=\sum_{j=1}^{M}m_j\,h_{x_0}(\rho_j).
\]
由附录定理 \ref{thm:app-comoving-horizon-localization} 的闭式表达，\(h_{x_0}(\rho)\) 是以 \((\gamma,\delta)\) 为中心/宽度参数的洛伦兹窗：
\[
\boxed{\;
h_{x_0}(\rho_j)=\frac{4\delta_j}{(\gamma_j-x_0)^2+(1+\delta_j)^2}\;}
\]
从而 \(H\) 为有限 Cauchy/Lorentz 核混合。

\begin{theorem}[扫描剖面的 Fourier--Laplace 指纹：有限指数混合与可辨识性]\label{thm:cdim-comoving-horizon-scan-fourier-inversion}
在上述设定下：
\begin{enumerate}
  \item \(H\in L^1(\RR)\cap C^\omega(\RR)\)；
  \item 其 Fourier 变换具有显式谱分解：对任意 \(\xi\in\RR\)，令
  \[
  \widehat H(\xi):=\int_{\RR}e^{-\mathrm{i}\xi x_0}H(x_0)\,\dd x_0,
  \]
  则
  \[
  \boxed{\;
  \widehat H(\xi)
  =4\pi\sum_{j=1}^{M}m_j\,\frac{\delta_j}{1+\delta_j}\,e^{-(1+\delta_j)|\xi|}\,e^{-\mathrm{i}\gamma_j\xi}\; } ;
  \]
  \item 若两组有限离线零点多重集诱导同一剖面 \(H\)（只需在任意非平凡开区间上一致），则两组多重集 \(\{(\gamma_j,\delta_j,m_j)\}\) 按重数一致。
\end{enumerate}
\end{theorem}
\begin{proof}
该定理即定理 \ref{thm:xi-comoving-horizon-scan-fourier-inversion} 在本章符号下的直接译文。证毕。
\end{proof}
该结论把二维参数 \((\gamma,\delta)\) 压缩为一维剖面 \(x_0\mapsto H(x_0)\) 的指数谱：\(\gamma\) 以平移协变的相位载体 \(e^{-\mathrm{i}\gamma\xi}\) 出现，而 \(\delta\) 以 \(e^{-(1+\delta)|\xi|}\) 的 Laplace 衰减出现；从而在有限缺陷口径下得到完全可逆的层析编码。

\begin{theorem}[最近缺陷层的 Fourier--Laplace 首层提取]\label{thm:cdim-comoving-horizon-scan-first-layer-extraction}
在定理 \ref{thm:cdim-comoving-horizon-scan-fourier-inversion} 的设定下，把不同深度值按升序排列为
\[
0<\delta^{(1)}<\delta^{(2)}<\cdots<\delta^{(r)},
\]
并定义首层指标集
\[
I_1:=\{j:\delta_j=\delta^{(1)}\}.
\]
若 \(r=1\)，约定 \(\delta^{(2)}:=+\infty\)，从而相应误差项视为 \(0\)。
则当 \(|\xi|\to\infty\) 时成立精确首层分离
\[
\boxed{\;
e^{(1+\delta^{(1)})|\xi|}\widehat H(\xi)
=
G_1(\xi)+
O\!\bigl(e^{-(\delta^{(2)}-\delta^{(1)})|\xi|}\bigr)\;}
\]
其中
\[
G_1(\xi):=
4\pi\sum_{j\in I_1}m_j\frac{\delta^{(1)}}{1+\delta^{(1)}}e^{-\mathrm{i}\gamma_j\xi}
\]
为有限三角多项式；若 \(r=1\)，则上式中的误差项按约定为 \(0\)。
特别地，最近临界层 \(\delta^{(1)}\) 及其全部横向位置 \(\{\gamma_j\}_{j\in I_1}\) 与重数 \(\{m_j\}_{j\in I_1}\) 可由高频尾部唯一恢复。
\end{theorem}
\begin{proof}
由定理 \ref{thm:cdim-comoving-horizon-scan-fourier-inversion} 的显式公式，按不同深度分组得
\[
\widehat H(\xi)
=
4\pi\sum_{\nu=1}^{r}
e^{-(1+\delta^{(\nu)})|\xi|}
\sum_{j:\delta_j=\delta^{(\nu)}}
m_j\frac{\delta^{(\nu)}}{1+\delta^{(\nu)}}e^{-\mathrm{i}\gamma_j\xi}.
\]
提出最小深度层 \(\delta^{(1)}\)，得到
\[
e^{(1+\delta^{(1)})|\xi|}\widehat H(\xi)
=
G_1(\xi)+
\sum_{\nu=2}^{r}
e^{-(\delta^{(\nu)}-\delta^{(1)})|\xi|}
G_\nu(\xi),
\]
其中各 \(G_\nu\) 仍为有限三角多项式，因而一致有界。若 \(r\ge 2\)，余项由最小正差
\(\delta^{(2)}-\delta^{(1)}\) 控制，遂得所示估计；若 \(r=1\)，则根本不存在余项。
\par
再证唯一恢复性。把 \(G_1\) 中具有相同横向位置的项合并，可写成
\[
G_1(\xi)=\sum_{\ell=1}^{r_1} c_\ell e^{-\mathrm{i}\eta_\ell\xi},
\qquad
c_\ell>0,
\]
其中 \(\eta_\ell\) 为首层不同 \(\gamma_j\) 的集合。有限频率指数函数族在 \(\RR\) 上线性无关，故该表示唯一；而
\[
c_\ell=
4\pi\frac{\delta^{(1)}}{1+\delta^{(1)}}
\sum_{j\in I_1:\,\gamma_j=\eta_\ell}m_j
\]
又唯一确定对应的总重数。于是首层深度、横向位置与重数全部由高频尾部唯一恢复。证毕。
\end{proof}
这说明共动视界扫描的高频尾部并非只给出“某种可辨识性”，而是带有天然的层级结构：最靠近临界线的一层总是先被读出，其后的各层只以更快衰减项出现。

\subsection{二维 Fourier--Laplace 全息层析：边界指纹唯一恢复缺陷测度}\label{subsec:cdim-fourier-laplace-holographic-tomography}
将有限多重集推广为 \(\RR\times(0,\infty)\) 上有限正测度 \(\nu\)。
定义边界密度
\[
H_\nu(x):=\int_{\RR\times(0,\infty)}\frac{4\delta}{(x-\gamma)^2+(1+\delta)^2}\,\dd\nu(\gamma,\delta),
\qquad
H_{\nu,t}:=P_t*H_\nu,
\]
其中 \(P_t(x)=\frac{1}{\pi}\frac{t}{x^2+t^2}\) 为 Poisson 核。
记 \(\widehat H_{\nu,t}(\xi)\) 为 Fourier 变换，令 \(s:=|\xi|\) 并定义二参数指纹
\[
F(\xi,s):=e^{(t+1)s}\widehat H_{\nu,t}(\xi).
\]

\begin{proposition}[二维指纹的张量化分解与单射性]\label{prop:cdim-defect-measure-fourier-laplace-holographic}
在上述设定下成立严格分解
\[
\boxed{\;
F(\xi,s)=4\pi\int_{\RR\times(0,\infty)}\frac{\delta}{1+\delta}\,e^{-s\delta}\,e^{-\mathrm{i}\gamma\xi}\,\dd\nu(\gamma,\delta)\; } ,
\]
从而 \(F\) 是 \(\gamma\) 的 Fourier 变换与 \(\delta\) 的 Laplace 变换之张量化组合；特别地映射 \(\nu\mapsto F\) 在有限测度类上为单射，因而缺陷测度 \(\nu\) 可由边界可见指纹 \(F\) 的取值唯一恢复。
\end{proposition}
\begin{proof}
该命题即结论 \ref{con:xi-terminal-defect-measure-fourier-laplace-holographic} 的直接表述。证毕。
\end{proof}

\begin{theorem}[有限缺陷测度的递归层剥离唯一性]\label{thm:cdim-defect-measure-recursive-layer-peeling}
在命题 \ref{prop:cdim-defect-measure-fourier-laplace-holographic} 的设定下，若
\[
\nu=\sum_{j=1}^{\kappa} m_j\,\delta_{(\gamma_j,\delta_j)}
\qquad
(m_j\in\ZZ_{\ge 1})
\]
为有限原子缺陷测度，并把不同深度值按升序排列为
\[
0<\delta^{(1)}<\delta^{(2)}<\cdots<\delta^{(r)},
\]
并约定 \(\delta^{(r+1)}:=+\infty\)。
则 \(\nu\) 可由 \(\widehat H_\nu(\xi)\) 通过有限步递归剥离唯一恢复。更精确地，若前 \(s-1\) 层已恢复并定义
\[
\nu^{(<s)}:=
\sum_{\nu'=1}^{s-1}\ \sum_{j:\delta_j=\delta^{(\nu')}}m_j\,\delta_{(\gamma_j,\delta_j)},
\]
以及残差 Fourier 变换
\[
\widehat H^{(s)}(\xi):=\widehat H_\nu(\xi)-\widehat H_{\nu^{(<s)}}(\xi),
\]
则对 \(1\le s\le r\) 有
\[
e^{(1+\delta^{(s)})|\xi|}\widehat H^{(s)}(\xi)
=
G_s(\xi)+
O\!\bigl(e^{-(\delta^{(s+1)}-\delta^{(s)})|\xi|}\bigr),
\]
其中
\[
G_s(\xi)=
4\pi\sum_{j:\delta_j=\delta^{(s)}}
m_j\frac{\delta^{(s)}}{1+\delta^{(s)}}e^{-\mathrm{i}\gamma_j\xi},
\]
特别地，当 \(s=r\) 时上式中的误差项自动视为 \(0\)。因此经过 \(r\) 步递归后，整个有限测度 \(\nu\) 被唯一确定。
\end{theorem}
\begin{proof}
由 \(H_\nu\) 的定义与命题 \ref{prop:cdim-defect-measure-fourier-laplace-holographic} 的 \(t=0\) 口径，
\[
\widehat H_\nu(\xi)
=
4\pi\sum_{\nu'=1}^{r}
e^{-(1+\delta^{(\nu')})|\xi|}
\sum_{j:\delta_j=\delta^{(\nu')}}
m_j\frac{\delta^{(\nu')}}{1+\delta^{(\nu')}}e^{-\mathrm{i}\gamma_j\xi}.
\]
当 \(s=1\) 时，这正是定理 \ref{thm:cdim-comoving-horizon-scan-first-layer-extraction} 的情形。设前 \(s-1\) 层已被剥离，则由构造有
\[
\widehat H^{(s)}(\xi)
=
4\pi\sum_{\nu'=s}^{r}
e^{-(1+\delta^{(\nu')})|\xi|}
\sum_{j:\delta_j=\delta^{(\nu')}}
m_j\frac{\delta^{(\nu')}}{1+\delta^{(\nu')}}e^{-\mathrm{i}\gamma_j\xi},
\]
这与首层提取定理完全同型，只需把最小深度替换为 \(\delta^{(s)}\)。于是
\[
e^{(1+\delta^{(s)})|\xi|}\widehat H^{(s)}(\xi)
=
G_s(\xi)+\text{更快衰减项},
\]
从而第 \(s\) 层可唯一恢复。对 \(s=1,\dots,r\) 归纳，即得整个有限原子测度 \(\nu\) 的唯一恢复。证毕。
\end{proof}
因此，二维 Fourier--Laplace 单射性不仅是静态存在性命题，而且给出一个按深度递归剥离的可执行恢复程序。

\subsection{低秩缺陷的有限采样可恢复性：Hankel--Prony 双向分解}\label{subsec:cdim-hankel-prony-reconstruction}
在命题 \ref{prop:cdim-defect-measure-fourier-laplace-holographic} 的设定下，若 \(\nu\) 为 \(\kappa\) 原子测度
\[
\nu=\sum_{j=1}^{\kappa}a_j\,\delta_{(\gamma_j,\delta_j)}
\qquad(\kappa\ge 1,\ a_j>0,\ (\gamma_j,\delta_j)\ \text{两两不同}),
\]
则二维指纹化为显式有限指数和
\[
F(\xi,s)=4\pi\sum_{j=1}^{\kappa}a_j\,\frac{\delta_j}{1+\delta_j}\,e^{-s\delta_j}\,e^{-\mathrm{i}\gamma_j\xi}.
\]

\begin{proposition}[有限采样的低秩可逆性与参数分离]\label{prop:cdim-atomic-defect-hankel-prony}
固定任意 \(\xi\neq 0\) 与步长 \(\Delta s>0\)，对任意起点 \(s_0>0\) 定义等距采样序列
\(
y_n:=F(\xi,s_0+n\Delta s)
\)。
则 \((y_n)\) 为 \(\kappa\) 项指数和
\[
y_n=\sum_{j=1}^{\kappa}c_j(\xi)\,\lambda_j^{n},
\qquad
\lambda_j=e^{-\Delta s\,\delta_j},
\qquad
c_j(\xi)=4\pi a_j\frac{\delta_j}{1+\delta_j}\,e^{-s_0\delta_j}e^{-\mathrm{i}\gamma_j\xi}.
\]
因此在有限缺陷/低秩口径下，\(\{\delta_j\}\) 可由 \(\{\lambda_j\}\) 的 Prony/Hankel（矩阵 pencil）指纹恢复，而 \(\{\gamma_j\}\) 可由 \(c_j(\xi)\) 的相位项 \(e^{-\mathrm{i}\gamma_j\xi}\) 恢复；并由二维 Fourier--Laplace 的单射性保证整体可辨识。
\end{proposition}
\begin{proof}
见结论 \ref{con:xi-terminal-atomic-defect-hankel-prony}。证毕。
\end{proof}

\begin{theorem}[单频 \texorpdfstring{$2\kappa$}{2kappa} 样本的精确深度恢复]\label{thm:cdim-atomic-defect-prony-2kappa-recovery}
在本小节设定下，固定任意 \(\xi\neq 0\)、步长 \(\Delta s>0\) 与起点 \(s_0\in\RR\)。若所有深度 \(\delta_1,\dots,\delta_\kappa\) 两两不同，并记
\[
y_n:=F(\xi,s_0+n\Delta s)\qquad(n\ge 0),
\]
则
\[
y_n=\sum_{j=1}^{\kappa} c_j(\xi)\lambda_j^n,
\qquad
\lambda_j:=e^{-\Delta s\,\delta_j},
\qquad
c_j(\xi):=4\pi a_j\frac{\delta_j}{1+\delta_j}e^{-s_0\delta_j}e^{-\mathrm{i}\gamma_j\xi}.
\]
在上述深度互异条件下，仅由前 \(2\kappa\) 个样本
\[
y_0,y_1,\dots,y_{2\kappa-1}
\]
即可唯一恢复全部深度 \(\delta_j\) 与复振幅 \(c_j(\xi)\)。
\leanverified{paper\_cdim\_atomic\_defect\_prony\_2kappa\_recovery}
\end{theorem}
\begin{proof}
由上式，\((y_n)\) 是 \(\kappa\) 项指数和。定义首一湮灭多项式
\[
Q(z)=\prod_{j=1}^{\kappa}(z-\lambda_j)
=
z^\kappa+q_{\kappa-1}z^{\kappa-1}+\cdots+q_0.
\]
则 \((y_n)\) 满足长度 \(\kappa\) 的线性递推
\[
y_{n+\kappa}+q_{\kappa-1}y_{n+\kappa-1}+\cdots+q_0y_n=0
\qquad(n\ge 0).
\]
取 \(n=0,\dots,\kappa-1\)，得到关于 \(q_0,\dots,q_{\kappa-1}\) 的 \(\kappa\) 个线性方程，其系数矩阵为 Hankel 矩阵
\[
\mathcal H_\kappa=(y_{i+j})_{0\le i,j\le \kappa-1}.
\]
由于 \(\lambda_j\) 两两不同且 \(c_j(\xi)\neq 0\)，有标准 Vandermonde 分解
\[
\mathcal H_\kappa=
V\,\diag(c_1(\xi),\dots,c_\kappa(\xi))\,V^{\mathsf T},
\qquad
V_{ij}=\lambda_j^i,
\]
从而 \(\mathcal H_\kappa\) 可逆。故湮灭多项式 \(Q\) 的系数唯一，进而其根 \(\lambda_j\) 唯一确定。于是
\[
\delta_j=-\frac{1}{\Delta s}\log \lambda_j
\]
唯一恢复全部深度。
\par
最后，把 \(\lambda_j\) 代回前 \(\kappa\) 个样本方程
\[
y_n=\sum_{j=1}^{\kappa} c_j(\xi)\lambda_j^n,
\qquad
0\le n\le \kappa-1,
\]
得到 Vandermonde 线性系统。因 \(\lambda_j\) 两两不同，故其解唯一，从而全部复振幅 \(c_j(\xi)\) 也被唯一恢复。证毕。
\end{proof}
这把低秩可逆性推进为一个精确样本复杂度结论：在深度互异的 generic 情形下，单频 \(2\kappa\) 样本恰足以恢复全部深度谱。

\begin{theorem}[双频 \texorpdfstring{$4\kappa$}{4kappa} 样本的二维参数精确恢复]\label{thm:cdim-atomic-defect-two-frequency-4kappa-recovery}
在定理 \ref{thm:cdim-atomic-defect-prony-2kappa-recovery} 的假设下，再取另一频率 \(\xi'\neq \xi\)，并假设所有横向参数满足先验窗口约束
\[
\gamma_j\in I\subset\RR,
\qquad
|I|<\frac{2\pi}{|\xi-\xi'|}.
\]
则仅由两组各 \(2\kappa\) 个样本
\[
\{F(\xi,s_0+n\Delta s)\}_{n=0}^{2\kappa-1},
\qquad
\{F(\xi',s_0+n\Delta s)\}_{n=0}^{2\kappa-1},
\]
即可唯一恢复整个有限原子族
\[
\{(\gamma_j,\delta_j,a_j)\}_{j=1}^{\kappa}.
\]
若进一步处在零点重数口径 \(a_j=m_j\in\ZZ_{\ge 1}\) 下，则全部 \((\gamma_j,\delta_j,m_j)\) 亦唯一恢复。
\end{theorem}
\begin{proof}
由定理 \ref{thm:cdim-atomic-defect-prony-2kappa-recovery}，第一组样本唯一恢复 \(\{\delta_j,c_j(\xi)\}\)，第二组样本唯一恢复 \(\{\delta_j,c_j(\xi')\}\)。由于深度两两不同，可据 \(\delta_j\) 对两组结果一一配对。对每个 \(j\)，有
\[
\frac{c_j(\xi)}{c_j(\xi')}
=
e^{-\mathrm{i}\gamma_j(\xi-\xi')}.
\]
于是
\[
\gamma_j
\equiv
-\frac{1}{\xi-\xi'}\Arg\!\left(\frac{c_j(\xi)}{c_j(\xi')}\right)
\pmod{\frac{2\pi}{|\xi-\xi'|}}.
\]
由先验窗口约束 \(|I|<2\pi/|\xi-\xi'|\)，模周期歧义被唯一解除，从而 \(\gamma_j\) 唯一确定。
\par
再由
\[
|c_j(\xi)|=4\pi a_j\frac{\delta_j}{1+\delta_j}e^{-s_0\delta_j},
\]
得
\[
a_j=
\frac{|c_j(\xi)|}{4\pi}\cdot \frac{1+\delta_j}{\delta_j}\,e^{s_0\delta_j},
\]
故全部权重 \(a_j\) 唯一恢复。若 \(a_j=m_j\in\ZZ_{\ge 1}\) 表示零点重数，则同式唯一确定相应重数。证毕。
\end{proof}
因此，在最自然的 generic 条件下，二维离线零点参数的恢复不再依赖无限采样，而可以压缩到两频 \(4\kappa\) 个样本的精确口径。

\subsection{扫描剖面的有限秩核化：从边界标量到 \texorpdfstring{$\kappa$}{kappa} 维再现核 Hilbert 空间}\label{subsec:cdim-scan-profile-finite-rank-rkhs}
对有限缺陷剖面 \(H\)（合并同参原子后可写为 \(\kappa\) 项洛伦兹混合）
\[
H(x)=\sum_{k=1}^{\kappa}\frac{w_k}{(x-\gamma_k)^2+(1+\delta_k)^2},
\qquad
w_k>0,\ \delta_k>0,
\]
引入复骨架点
\(
p_k:=\gamma_k-\mathrm{i}(1+\delta_k)\in\CC_{-}
\)
与特征映射
\(
v:\RR\to\CC^\kappa
\),
\(
v(x)_k:=\frac{\sqrt{w_k}}{x-p_k}
\)。

\begin{theorem}[有限秩核的对角核化与 \texorpdfstring{$\kappa$}{kappa} 维 RKHS]\label{thm:cdim-basepoint-scan-finite-rank-rkhs}
定义核
\[
K(x,y):=\langle v(x),v(y)\rangle_{\CC^\kappa}
=\sum_{k=1}^{\kappa}\frac{w_k}{(x-p_k)(y-\overline{p_k})}
=\sum_{k=1}^{\kappa}\frac{w_k}{(x-\gamma_k+\mathrm{i}(1+\delta_k))(y-\gamma_k-\mathrm{i}(1+\delta_k))}.
\]
则 \(K\) 为正定核并满足对角恒等式
\[
\boxed{\;K(x,x)=H(x)\qquad(x\in\RR).\;}
\]
从而扫描剖面可被精确核化为一个有限秩核的对角；并且
\(\mathcal{H}:=\overline{\Span}\{v(x):x\in\RR\}\subseteq\CC^\kappa\)
在自然内积下为 \(\kappa\)-维 Hilbert 空间（等价地：\(K\) 为一维参数生成的 \(\kappa\)-维 RKHS 的再现核）。
\end{theorem}
\begin{proof}
该定理即定理 \ref{thm:xi-basepoint-scan-finite-rank-rkhs} 在本章口径下的直接译文。证毕。
\end{proof}
该核化把“边界标量剖面”提升为“有限秩核几何”，从而把层析反演的稳定性与数值审计接口严格对齐到有限秩核在采样点上的 Gram 数据恢复问题；这与 Hankel--Prony/矩阵 pencil/ESPRIT 的标准反演口径兼容。

\subsection{位复杂度屏障与共动前缀收益：二次压缩的制度成本与可审计回收}\label{subsec:cdim-bit-budget-barrier-comoving-prefix-gain}
\subsubsection{端点二次压缩与高度窗上的最坏情形位预算}
固定 Cayley 图表下，离临界线偏移在端点层余量 \(h(\rho)=1-|w_\rho|^2\) 中发生二次压缩（定理 \ref{thm:xi-offcritical-quadratic-radial-compression}）。
在高度窗 \(|\gamma|\le T\) 且偏移下界 \(\delta\ge\delta_0\) 的最坏情形中，端点余量满足显式下界（定理 \ref{thm:xi-offcritical-endpoint-resolution-lower-bound}）
\[
1-|w_\rho|^2\ \ge\ \frac{4\delta_0}{T^2+(1+\delta_0)^2},
\]
从而若以 \(2^{-p}\) 作为端点余量的最小可分辨尺度，则必有位预算下界
\[
p\ \ge\ \log_2\!\Bigl(\frac{T^2+(1+\delta_0)^2}{4\delta_0}\Bigr)
=2\log_2T+\log_2\!\frac{1}{\delta_0}+O(1)\qquad(T\to\infty),
\]
并可在精度势口径下等价写成 \(\mathcal{P}(T)=\log\pi+\log(1+T^2)\) 的显式形态（式 \eqref{eq:xi-offcritical-endpoint-bit-precision-potential-form}）。

\subsubsection{共动前缀的常数预算与制度成本}
若允许在对象层外置可寻址前缀 \(x_0\) 并取共动局部化 \(x_0=\gamma\)，则
\[
h_{\gamma}(\rho)=1-|w_{\rho,\gamma}|^2=\frac{4\delta}{(1+\delta)^2}\ \ge\ \frac{4\delta_0}{(1+\delta_0)^2}\qquad(\delta\ge\delta_0),
\]
从而端点余量阈值可在高度无关的常数位预算下证书化（推论 \ref{cor:xi-comoving-prefix-bit-budget-gain}）。
需要强调的是：该“回收”并不消失，而是被制度性转移为地址前缀预算；若前缀未给定，则相应评价在对象层不定义并读出为 \(\mathsf{NULL}\)（注 \ref{rem:xi-comoving-horizon-localization-address-null}）。

\subsubsection{紧化读出的径向位复杂度下界}
相位紧化把动态范围强制压入末位精度（式 \eqref{eq:app-precision-amplification}），并在“仅允许输出 dyadic 括号”的制度中导出硬下界：
若仅通过 \(u=\upsilon(x)=\pi^{-1}\arctan(x)\) 的 \(n\)-比特 dyadic 近似来证书化 \(x\) 且要求离线复核得到绝对误差 \(\le\varepsilon\) 的单值可核验读出，则必有（定理 \ref{thm:xi-radial-information-projection-lower-bound}）
\[
\boxed{\;
n\ \ge\ \log_2\!\Bigl(\frac{1}{\varepsilon}\Bigr)+\log_2(1+x^2)+\log_2\pi+O(1)\; }.
\]
将其应用到完成化接口的径向参数 \(x=\log|u_L(s)|=(2\Re(s)-1)\log L\) 可见：当 \(\Re(s)\neq\tfrac12\) 且允许 \(L\to\infty\) 时，证书位数不可能对 \(L\) 一致有界，必须随 \(\log_2(1+x^2)\) 增长（同上定理）。

\subsection{正交切片定理：可见模长分量、盘内零点分量与端点载荷分量}\label{subsec:cdim-orthogonal-slice-theorem}
在 unitary slice 可见域约束（\(X_{\mathrm{vis}}=\mathsf{Phase}\)）与 \(\mathbf{PW}\) 闭包类型检查下，偏离临界切片的信息在可审计意义下被强制分解为三类互不替代的正交通道：
\begin{enumerate}
  \item \textbf{协议正交（径向分量）：}任何本质上需要 \(|u|\neq 1\) 的连续径向自由度才能表达的对象，在纯相位接口中不可类型化；若拒绝外置唯一径向寄存器 \(\varrho=\log|u|\)，则读出按协议退化为 \(\mathsf{NULL}\)（定理 \ref{thm:xi-pw-no-continuous-hair}、推论 \ref{cor:xi-unique-continuous-transverse-register}，并见 \S\ref{subsec:cdim-offline-zero-claim-interface}）。
  \item \textbf{因子分解正交（盘内零点分量）：}在 Hardy/Smirnov 因子分解口径下，边界模长仅决定外因子，而盘内零点与奇异支撑完全由内因子承载；因此“离切片点状事件”的零点数据属于内因子分量，对边界模长审计结构性不可见（附录 \S\ref{subsec:unit-disk-inner-outer}）。
  \item \textbf{端点载荷正交（端点分量）：}端点邻域的伸缩半群作用分解为有限离散残差与连续精度势负载两条互不耦合的账本轴（\S\ref{subsec:cdim-endpoint-orthogonal-load}），从而端点载荷在证书闭包中对应独立记录通道，不能由提升半径采样或地址深度替代（并见 \S\ref{subsec:cdim-orthogonal-threshold-nonsubstitutable}）。
\end{enumerate}
因此，“离开 \(\Re(s)=\tfrac12\)”若要在闭环内保持可核验，只能沿上述正交不变量与外置预算轴出现；拒绝外置其中任一支撑轴，将在类型层必然触发 \(\mathsf{NULL}\) 或其可审计失败见证（定理 \ref{thm:xi-null-complete-trichotomy-offline}）。

\endinput
